Necessarily $m>n$; put $D=m-n$. The condition is $n+N\mid D$, so the $20$ admissible values correspond to divisors of $D$ greater than $n$. In particular $\tau(D)\ge21$.

The only positive integer below $368$ with at least $21$ divisors is $360$. To verify this without a long scan, write prime exponents in decreasing order. One prime needs exponent at least $20$. For two primes the least possible values are at least $2^6 3^2=576$ (the other exponent patterns are larger). For three primes, the only pattern that can be below $368$ and have at least $21$ divisors is $(3,2,1)$, whose least value is $2^3 3^2 5=360$; changing a prime makes it larger than $368$. Pattern $(5,1,1)$ is already $480$, and patterns at least $(2,2,2)$ are at least $900$. Four primes need a repeated prime, giving at least $420$; five primes give at least $2310$.

For $D=360$, $\tau(D)=24$. Exactly four divisors must be at most $n$. They are $1,2,3,4$, so the smallest such $n$ is $4$. Taking $m=364$ works and gives $m+n=368$. Any $D\ge368$ gives $m+n=D+2n>368$. Thus the minimum is $\boxed{368}$.
